\section{Introduction}\label{sec:intro}
In passenger railway service planning, the train timetabling stage determines the arrival and departure times of trains and thus how and when trains use the limited infrastructure of a line \citep{Cacchiani2012}. A timetable is only executable if train movements between stations are compatible with the tracks and routes available inside stations. Section running times and minimum headways alone do not ensure this: two trains may request the same arrival--departure track at overlapping times, or use conflicting routes through the station throat area, even when all section headways are respected \citep{Lusby2011survey,Wang2023}. These station-level interactions are usually resolved after timetabling, in the train platforming problem \citep{Caprara2011}. When the timetable is tight, however, a sequential approach can leave conflicts that cannot be removed without changing the timetable itself.

Optimizing stopping patterns together with the timetable adds a further layer of dependency. A stop changes the dwell time, the running times on the adjacent sections, the possibility of being overtaken, and the set of admissible station routings. Skipping a stop can therefore remove a local station conflict, but it may violate a station service frequency requirement; another train may then have to stop instead, which changes its own times and station routing. Even if no stop is changed, shifting an existing stop across the boundary of a counting time window changes the number of services counted in that window. A repair that only reconsiders the trains involved in a conflict may therefore lack the decisions needed to restore feasibility. Stop planning, timetabling, and platforming are interdependent and should be optimized together \citep{Yang2016,Dong2020,Xu2021}.

From an industrial engineering perspective, the problem is a large scheduling problem with shared, disjunctive resources (section block, station tracks, and routes) and with service-level constraints that couple many jobs (trains) through counting rules. It is natural to represent such a problem on an event-activity network, in which arrival and departure events are linked by running, dwell, and headway activities \citep{Serafini1989,Liebchen2007}, and in which the resource conflicts between trains appear as disjunctive pairs of activities, as in the alternative graph of job-shop scheduling \citep{Mascis2002,DAriano2007}. In our setting the network is not fixed: which station activities exist depends on the stopping and routing decisions, and which stops are counted depends on the event times.

This paper studies the integrated train timetabling, stop planning, and platforming problem for prescribed train services on a railway line. We determine stopping patterns, arrival and departure times, train orders, and station routings, i.e., an arrival--departure track (or main track) together with compatible receiving and departure routes for every train at every station. Train stop-count constraints, station service frequency constraints, maximum travel times, and maintenance windows restrict these decisions. The primary objective is to minimize the timetable makespan, i.e., the time between the earliest origin departure and the latest destination arrival, so that the prescribed train services are accommodated within the shortest operating period; deviations from preferred departure and arrival windows and running time supplements are secondary objectives. All prescribed trains must be operated, so the makespan is reduced through coordinated scheduling rather than cancellations.

The full model is difficult to solve directly. Its station routing and pairwise order variables grow quickly with the number of train operations at each station; in our large-scale case the full mixed-integer linear program (MILP) has about 14 million constraints and 0.7 million binary variables. We therefore develop a conflict-driven adaptive large neighborhood search (CD-ALNS), a matheuristic in which destroy operators release a subset of decisions selected from the conflict structure of the current solution, and a restricted MILP repairs them while all constraints linking the released and the fixed decisions remain active.

The main contributions of this paper are as follows.
\begin{enumerate}
\item We represent the integrated problem on an event-activity network with disjunctive and conditional activities, and formulate it as a MILP with continuous event times. The formulation couples train stop-count and station service frequency constraints, which are evaluated on closed time windows, with section headways, overtaking, track occupation, and route headways. We show that, for a fixed selection of stops, orders, and station routings, the minimum makespan equals the length of a longest path in the selected network, which identifies the decisions that must change to reduce the makespan.
\item We develop a CD-ALNS with three phases and nine destroy operators. A macroscopic phase restores the feasibility of train movements and service frequencies; a joint feasibility phase resolves station conflicts; and a joint quality phase reduces the makespan by releasing the decisions on critical activities. The joint neighborhoods include the compensating stops and event shifts needed to keep every affected service count within its bounds, so that station conflicts and service frequency requirements are resolved in the same restricted MILP.
\item We evaluate the approach on three cases with up to 29 stations and 101 trains, against direct MILP solution and an adapted logic-based Benders decomposition (LBBD), and analyse the convergence, the resulting timetables and station track usage, and the contribution of the macroscopic phase.
\end{enumerate}

The remainder of the paper is organized as follows. \Cref{sec:lit} reviews the related literature. \Cref{sec:problem} describes the problem and its event-activity network representation, and \cref{sec:model} presents the mathematical formulation and its structural properties. \Cref{sec:alg} introduces the CD-ALNS. \Cref{sec:exp} reports the computational experiments, \cref{sec:insights} discusses managerial implications, and \cref{sec:concl} concludes the paper.
